%%========================================================================
%% WAJIB DIISI SESUAI DATA YANG BENAR DAN TIDAK BOLEH ADA YANG TERLEWAT
%%========================================================================

\providecommand{\judulid}{RANCANG BANGUN APLIKASI POINT OF SALE DENGAN FITUR PREDIKSI PENJUALAN HARIAN MENGGUNAKAN METODE LONG SHORT-TERM MEMORY (LSTM)}
\providecommand{\judulen}{DESIGN AND DEVELOPMENT OF A POINT OF SALE APPLICATION WITH DAILY SALES PREDICTION FEATURE USING THE LONG SHORT-TERM MEMORY (LSTM) METHOD}
\providecommand{\penulis}{MUHAMMAD KHOLIS}
\providecommand{\nim}{2022573010098}

\providecommand{\tipe}{Skripsi}
\providecommand{\type}{Undergraduate Thesis}
\providecommand{\prodi}{Teknik Informatika}
\providecommand{\fakultas}{Teknologi Informasi dan Komputer}
\providecommand{\universitas}{Politeknik Negeri Lhokseumawe}

\providecommand{\tglpersetujuan}{28 Mei 2026}
\providecommand{\tglpernyataan}{28 Mei 2026}
\providecommand{\tglpengesahan}{28 Mei 2026}
\providecommand{\tahun}{2026}

% Pembimbing
\providecommand{\pembimbingI}{Zulfan Khairil Simbolon, S.T., M.Eng.}
\providecommand{\NIPpembimbingI}{19690902 199303 1 004}
\providecommand{\pembimbingII}{Azhar, S.T., M.T.}
\providecommand{\NIPpembimbingII}{19640830 199003 1 005}
\providecommand{\pembimbing}{Zulfan Khairil Simbolon, S.T., M.Eng.}
\providecommand{\NIPpembimbing}{19690902 199303 1 004}

% Ketua Jurusan & Kaprodi
\providecommand{\kajur}{Salahuddin, S.T., M.Cs}
\providecommand{\NIPkajur}{19740424 200212 1 001}
\providecommand{\koorprodi}{M. Khadafi, S.T., M.T}
\providecommand{\NIPkoorprodi}{19750718 200212 1 004}

% Penguji
\providecommand{\pengujiI}{Dr. Rahmad Hidayat, S.Kom., M.Cs}
\providecommand{\NIPpengujiI}{19830420 201212 1 003}
\providecommand{\pengujiII}{Muhammad Rizka, S.S.T., M.Kom}
\providecommand{\NIPpengujiII}{19881009 201504 1 001}
\providecommand{\pengujiIII}{Bismillah Baik}
\providecommand{\NIPpengujiIII}{-}
